%!TEX program = lualatex

\documentclass[portuguese, 11pt, a4paper, titlepage, oneside]{article}

\usepackage{geometry}
\usepackage[T1]{fontenc}
\usepackage{graphicx}
\usepackage{parskip}
\usepackage{amssymb}
\usepackage[portuguese]{babel}
\usepackage[framed, numbered]{matlab-prettifier}
\usepackage{subcaption}
\usepackage{wrapfig}
\usepackage{hyperref}

% Title page config
\title{
  \includegraphics[width=10cm]{assets/deti_horizontal_preto.pdf}
  \vskip 2em
  \huge Aspetos Profissionais e Sociais\\
  da Engenharia Informática
  \vskip 1em
	{\bfseries Trabalho individual 1}
}

\author{Rúben Garrido\\NMec: 107927}
\date{\today}
% Begin document
\begin{document}

% Title
\maketitle

\section{Utilização do ChatGPT pelos alunos e professores, no contexto de APSEI}
\subsection{Introdução}
A inteligência artificial baseada em \textit{Large Language Models} (neste caso, o ChatGPT) revolucionou a forma como os humanos pensam e interagem com a tecnologia. Todas as áreas foram influenciadas, embora a educação tenha tido um dos maiores destaques.

Esta mudança de paradigma deve-se sobretudo à utilização do ChatGPT (bem como de outros LLMs) pelos alunos, por forma a "facilitar" a curva de pensamento necessária para realizar um determinado trabalho. No entanto, o uso excessivo e irresponsável da inteligência artificial levou ao surgimento de respostas meramente criadas por esta, com imprecisões e padrões repetidos.

Aspetos Profissionais e Sociais da Engenharia Informática é uma Unidade Curricular que, ao contrário das restantes UCs do curso, se foca em questões éticas e de posicionamento dos engenheiros informáticos perante a sociedade. Esta visa promover a formação de profissionais conscientes do mundo ao seu redor a nível legislativo, financeiro e social, bem como das implicações da tecnologia no panorama atual.

O objetivo deste trabalho foca-se na perceção da usabilidade do ChatGPT no âmbito do processo de ensino e aprendizagem da unidade curricular APSEI, por forma a um melhor entendimento geral no que concerne ao processo de colaboração ou comprometimento dos processos de ensino.

\subsection{Ética e moral}
À primeira instância, o que me ocorre de imediato é "Não! Seria imoral, enganador...". Dependendo em que situação ou de que forma seria utilizado o ChatGPT, como tudo na sociedade existem pontos positivos e pontos negativos, e certamente que o uso desta ferramenta, na UC de APSEI poderia levantar questões éticas e imorais, uma vez que certos trabalhos apelam à criatividade de cada aluno, que, ao utilizarem a inteligência artificial, poderiam colocar em causa a hipótese ou não de um plágio, estando o aluno a pôr em questão a sua ética profissional.

Se, neste momento, para realizar este trabalho eu recorresse à IA por completo, estaria a reduzir a minha capacidade criativa e original de pensamento, e, moralmente, estaria a enganar o docente, o que eticamente, a nível de integridade académica não é correto. APSEI requer muita criatividade artística de pensamento "fora da caixa", o que, para alguns alunos, poderá ser a "pedra no sapato" do curso. Recorrendo ao ChatGPT, os alunos poderão ver ali o escape e, ao mesmo tempo, a solução para o problema que a maioria enfrenta da escrita. Compete aos alunos terem o discernimento de usarem esta ferramenta de forma útil mas responsável, quanto aos aspetos éticos da unidade curricular.

\subsection{Docente, alunos e disciplina}
Sendo que o aluno é avaliado pela sua capacidade de compreensão da matéria fornecida, é necessário compreender se o ChatGPT fará sentido no âmbito da disciplina. A meu ver, a utilização do ChatGPT pelos alunos nesta unidade curricular \textbf{deve ser permitida, e até expectável} por parte dos professores. No entanto, e como tudo, deve ser utilizada com moderação e com espírito crítico.

Na minha opinião, não há necessidade de incentivar o aluno a utilizar o ChatGPT, uma vez que a IA já começa a estar enraizada na sociedade, em particular na camada mais jovem. Deve-se sim é incentivar o aluno a usá-lo apenas como auxiliar de aprendizagem, para que não perca as competências cognitivas e críticas.

Se o aluno encontra dificuldades na compreensão ou até mesmo na exposição das suas ideias, esta ferramenta será obviamente uma mais-valia à sua formação, organização de raciocínio e argumentação. Tudo vai depender da sua capacidade de questionar o ChatGPT (através dos \textit{prompts} realizados), confrontar as informações dadas, pegar nelas, obter espírito crítico e ir mais além no pensamento do que aquilo que lhe é fornecido. Aqui, torna-se uma ferramenta útil e favorável, e é mais um acesso a uma via de informação, que positivamente serve para potenciar a aprendizagem. Compete aos alunos entender os limites desta ferramenta e tomar decisões informadas sobre quando e como utilizá-la de maneira apropriada.

Resumindo, do ponto de vista da aprendizagem, o aluno deve ser autorizado a usar o ChatGPT no âmbito da disciplina de APSEI, de forma cuidada e responsável, com o intuito de progredir na compreensão da matéria. Os alunos devem esperar que o ChatGPT os auxilie e não que lhes faça o trabalho na íntegra.

Em relação ao docente, o ChatGPT pode até ser benéfico, uma vez que pode gerar uma resposta imediata ao aluno. No entanto, um dos maiores desafios é a necessidade de analisar as respostas ou trabalhos dados pelo aluno, uma vez que podem ser apresentados de uma forma vaga e aleatória e a informação nem sempre é 100\% precisa.

Uma questão que se impõe é: até que ponto é que um docente pode complementar a avaliação de um aluno de forma justa, através do ChatGPT? Respondendo a esta questão, sou da opinião de que o docente não deve perguntar ao ChatGPT se um determinado trabalho está certo ou errado, visto que a IA se baseia em probabilidades e, consequentemente, o aluno pode ser avaliado de uma forma errada. Deve questionar pontos favoráveis e pontos negativos e, com base nessa informação e numa observação externa, formular uma opinião que resultará numa avaliação. Mais uma vez, a regra que deve ser aplicada aos alunos repete-se ao docente, com a ferramenta a apenas complementar e não substituir a competência humana.

O docente pode usar o ChatGPT para promover o interesse da matéria lecionada; pedir para os alunos retirarem informação básica para mais tarde ser discutida em grupo na aula, promove o pensamento crítico do aluno e intensifica a interação humana.

Em suma, o docente pode tirar partido também como ferramenta de apoio, nunca esquecendo que a relação docente-aluno será sempre prioritária, promovendo a socialização, empatia e interação humana entre ambas as partes.

A disciplina passa a ser, para o aluno que apresente dificuldades na matéria, algo menos maçudo e pesaroso, originando uma maior vontade de a frequentar, dando a facilidade de ajuda na compreensão da matéria, na criação do pensamento crítico e da análise.

Sendo que APSEI é uma disciplina de pensamento "fora da caixa", o aluno poderá recorrer ao ChatGPT para inicialmente se contextualizar, obter informações primárias e, posteriormente, transformá-las ou complementá-las, mobilizando os seus pensamentos e conhecimentos, por forma a encontrar a resposta ao pretendido. Nem questiono em relação à acessibilidade e à equidade, uma vez que, sendo APSEI uma UC do curso de Engenharia Informática, parte-se do princípio de que cada aluno tem acesso à Internet e ao ChatGPT, bem como a destreza para utilizar esta ferramenta.

A integração da inteligência artificial vem trazer algumas ideias novas, em situações onde nem sempre se consegue discernir por nós próprios. Esta permite simular alguns dilemas éticos, e inclusive olhar para um certo problema com uma visão diferente.

Para além disso, a utilização do ChatGPT pode ser benéfica \textbf{no que toca à correção linguística, gramatical e ortográfica}. Uma vez que APSEI é a única unidade curricular do curso que onde é necessário realizar trabalhos de escrita, é comum os alunos não se sentirem tão confortáveis a escrever artigos de opinião, devido à falta de prática.

Assim, e por um lado, o ChatGPT permite corrigir erros gramaticais e reescrever frases menos bem construídas. Por outro, quando utilizada com espírito crítico, pode ser um incentivo aos alunos de modo a que estes desenvolvam \textit{skills} linguísticas mais aprofundadas.
\pagebreak

\section{Deverão os professores gravar as aulas e tornas os vídeos disponíveis para os alunos mais tarde?}
As gravações das aulas e posteriormente o envio das mesmas em vídeo não garantem por si só a aprendizagem do aluno. São apenas um recurso ou complemento ao aluno, de forma a possibilitar o alcance dos seus objetivos em relação a disciplina. Neste sentido, o aluno ou por falta à aula, ou por distração, ou mesmo por incompreensão da mesma naquele dia, poderá rever a aula e tentar, andando para a frente e para trás no vídeo, compreender ou esclarecer as suas dúvidas.

As visualizações das aulas em vídeos são mais sensoriais e visuais, em contrapartida com as aulas presenciais, que são mais racionais e abstratas. Isto porque, no primeiro caso, o aluno lê e vê, podendo voltar atrás em caso de dúvida e absorver com mais calma a informação ao seu próprio ritmo. No segundo caso, o aluno escuta e na hora já raciocina sobre o que está ouvir, não dando tanto azo a que se possa parar a aula que o professor dá, numa situação em que existam muitos alunos a assistir à mesma. Isto porque seria impossível o professor esclarecer numa só aula 200 alunos ao ritmo de compreensão de cada um.

Neste sentido, as aulas gravadas são uma mais valia, pois o aluno pode chegar a casa, rever a aula já disponibilizada em vídeo, tentar perceber o que não compreendeu e, se mesmo assim tiver dificuldades, na aula seguinte pedir ajuda ao professor.

A gravação das aulas pode obrigar o decente a um esforço adicional, no que concerne a aquisição de novas competências. Ao mesmo tempo, as aulas gravadas podem também ter um efeito negativo, uma vez que o aluno que é organizado irá sempre que possível assistir às aulas presenciais, e manter a matéria atualizada. No entanto, aqueles estudante que é mais preguiçoso irá apenas se sujeitar aos vídeos, acumulando matéria até à altura dos exames.

Acredito que a maior parte dos estudantes utiliza as aulas gravadas nas disciplinas em que a matéria é de fácil aprendizagem ou contém muito material gráfico, ou até mesmo como complemento para a revisão da matéria na altura dos exames.

As aulas presenciais permitem interagir com o docente em caso de dúvidas, servindo as gravadas apenas para consolidação da matéria e complementação de informação.

Ao trabalhador estudante, estas são uma vantagem notória, podendo o acesso aos vídeos gravados mantê-los ao nível dos colegas, caso trabalhem em horários de aulas presenciais.

No entanto, as aulas gravadas apresentam desvantagens caso sejam leituras narradas de diapositivos, pois promovem uma menor interação aluno-docente. No caso de APSEI, sendo que as aulas se baseiam bastante na explicação da matéria por parte do professor (ao invés de serem à base de leitura de diapositivos), a gravação das aulas torna-se mais útil do que nas restantes UCs, já que nem sempre os \textit{slides} refletem todo o conteúdo lecionado.

As aulas práticas de APSEI, a meu ver, não fazem tanto sentido serem gravadas como as teóricas, porque, aquando da avaliação crítica do professor a cada aluno ou grupo, esta seria exposta caso fosse gravada. Nesse sentido, penso que cada avaliação é particular e, portanto, não deva ser exteriorizada para todos.

Os vídeos devem ser partilhados para todos os alunos sem exceção, cabendo a cada um deles ter consciência que o facto de os ter serve apenas para complementar e não substituir a aula presencial.

Convém ter em conta que a gravação das aulas e a sua posterior divulgação aos alunos pressupõe o consentimento de todos os que foram visualizados, sejam eles alunos ou docentes, no que diz respeito à proteção de dados pessoais e à livre circulação desses mesmos dados. Daí que, tanto a intervenção dos alunos, bem como o uso dos microfones que apanhem o áudio de todos tenha de ser bem considerado e questionado.

Num contexto geográfico, as aulas gravadas, a título da União Europeia, devem obedecer a certas regras de privacidade e proteção de dados. Num país como, por exemplo, os Estados Unidos da América, penso que, como não existe tanta exigência legislativa no que concerne a este tema e, portanto, a gravação das aulas é menos condicionada.

Um ponto muito importante prende-se com o facto de toda e qualquer palavra ou frase que, num determinado contexto, possa ser proferida, quer pelo professor, quer por qualquer aluno, para uma determinada pessoa, possa ferir suscetibilidade. Nesta situação, a gravação das aulas pode-se tornar negativa, sendo que uma pequena coisa pode dar origem a uma "grande bolha desnecessária". Uma medida para isto passaria por restringir o acesso a um determinado grupo de pessoas (isto é, o conjunto de alunos que frequentam APSEI), embora isso não impeça o surgimento de \textit{leaks} que comprometam a pessoa em questão numa área abrangida para além do grupo restrito.

Em suma, a gravação das aulas e consequente disponibilização dos vídeos, a meu ver, tem tudo para ser positivo nas aulas teóricas de APSEI, não fazendo tanto sentido nas aulas práticas. Provavelmente, poderá haver um esforço adicional para o docente, uma vez que para além de lecionar as aulas, ainda terá o trabalho de as gravar e disponibilizar. Há que ter em atenção as regras de privacidade e proteção de dados. Compete aos alunos usar as aulas como um benefício complementar aos estudos e não em detrimento da assistência presencial à aula. Em turmas com um elevado número de alunos, a gravação irá auxiliar a aprendizagem e compreensão da matéria lecionada.
\end{document}